\rok{Јануар 2024.}{B - Поправни тест другог теста}

Поправни тест из Алгоритама и структура података

Други практични тест одржан 19.01.2024.

\zadatak{1}{Завршавање имплементације BFS алгоритма}
Основна метода BFS алгоритма није потпуна. Допунити имплементацију ове методе која треба да омогући комплетно извршавање овог алгоритма.

\textbf{Додатне напомене за решавање задатка:}
\begin{itemize}
	\item Написати алгоритам.
	\item У Main методи креирати произвољан граф. Обезбедити начин за тестирање и проверу нове функционалности, односно позив BFS методе.
	\item У template-у је закоментарисана метода:
	\begin{Verbatim}[fontsize=\footnotesize, numbers=left, xleftmargin=10mm]
public void BFS(int startVertex)
	\end{Verbatim}
\end{itemize}



\zadatak{2}{Провера да ли је низ бројева палиндром коришћењем Hash табеле}
Написати алгоритам временске комплексности $O(n/2)$ који ће употребом Hash табеле проверити да ли је неки низ бројева палиндром. Палиндром је низ бројева који читан од почетка до краја и обрнуто изгледа идентично.

\textbf{Додатне напомене за решавање задатка:}
\begin{itemize}
	\item Обавезно је користити Hash табелу која је дата у шаблону.
	\item Може се креирати више Hash табела.
	\item Велика и мала слова посматрати различито.
	\item Имплементацију писати у оквиру закоментарисане методе:
	\begin{Verbatim}[fontsize=\footnotesize, numbers=left, xleftmargin=10mm]
public static bool IsPalindrome(int[] arr)
	\end{Verbatim}
\end{itemize}

\textbf{Примери:}
\begin{itemize}
	\item \textbf{Улаз 1:} \texttt{[1, 2, 3, 4, 5]}
	\item \textbf{Излаз 1:} false
	\item \textbf{Улаз 2:} \texttt{[1, 2, 3, 2, 1]}
	\item \textbf{Излаз 2:} true
\end{itemize}



\zadatak{3}{Исписивање чворова који су листови са парном вредношћу}
У оквиру стандардне структуре класе \texttt{BST} написати имплементацију методе која ће омогућити пребројавање чворова који се налазе на позицији листова, а имају парну вредност.

\textbf{Додатне напомене за решавање задатка:}
\begin{itemize}
	\item У оквиру класе \texttt{BST} креирати јавну методу са потписом:
	\begin{Verbatim}[fontsize=\footnotesize, numbers=left, xleftmargin=10mm]
public void countEvenLeaves()
	\end{Verbatim}
	\item Из претходно креиране јавне методе треба да се позива приватна метода са потписом:
	\begin{Verbatim}[fontsize=\footnotesize, numbers=left, xleftmargin=10mm]
private int countEvenLeavesAVL(Node curr)
	\end{Verbatim}
	\item Приватна метода треба да садржи имплементацију алгоритма којим ће се омогућити пролазак кроз стабло од корена ка листовима, исписивање чворова са парном вредношћу.
	\item Листом се класификује сваки чвор који нема ниједно дете.
	\item Алгоритам \textbf{ОБАВЕЗНО} писати у оквиру класе \texttt{BST}.
\end{itemize}

\textbf{Примери:}

\begin{center}
\begin{minipage}{\textwidth}
\centering
\begin{forest}
for tree={
	circle,
	draw,
	thick,
	fill=gray!20,
	minimum size=8mm,
	inner sep=1pt,
	s sep=8mm,
	l sep=12mm,
	edge={-, thick},
	font=\small
}
[8
	[4
		[2
			[, phantom]
			[3]
		]
		[5]
	]
	[9
		[, phantom]
		[11]
	]
]
\end{forest}

\vspace{0.3cm}
\textbf{Слика 1.} Пример BST-а.
\end{minipage}
\end{center}

\begin{itemize}
	\item \textbf{Улаз 1:} BST са слике 1
	\item \textbf{Излаз 1:} \texttt{/}
\end{itemize}

\begin{center}
\begin{minipage}{\textwidth}
\centering
\begin{forest}
for tree={
	circle,
	draw,
	thick,
	fill=gray!20,
	minimum size=8mm,
	inner sep=1pt,
	s sep=8mm,
	l sep=12mm,
	edge={-, thick},
	font=\small
}
[10
	[4
		[2
			[3]
			[, phantom]
		]
		[6
			[5]
			[9]
		]
	]
	[20
		[12
			[, phantom]
			[14]
		]
		[29
			[27]
			[, phantom]
		]
	]
]
\end{forest}

\vspace{0.3cm}
\textbf{Слика 2.} Пример BST-а.
\end{minipage}
\end{center}

\begin{itemize}
	\item \textbf{Улаз 2:} BST са слике 2
	\item \textbf{Излаз 2:} 14
\end{itemize}



\sablon{Шаблон поправног за други тест јануар 2024. група Б}{2023/Januar/GrupaB/AISP2023_PopravniTest2_GrupaB.linq}
